\begin{pastexam}[label={px:ntut-108-ci-03}]{北科大資工所 108 學年度 第 1 題（節錄 (5) 小題）}
（簡答）

(15 pts) Please answer the following questions concisely.
\begin{enumerate}[label=(\arabic*)]
\setcounter{enumi}{4}
\item (3 pts) Let A and B be two problems and we were able to establish the following fact: ``if we
could solve A in time O($T(n)$), then we could solve B in time O($n\log n+T(n)$).'' Suppose that B
now has a $\Omega$($n\log n$) time lower bound. Can we conclude that A also has a $\Omega$($n\log n$)
time lower bound?
\end{enumerate}
\end{pastexam}

\begin{pastsoln}
答案：不能。

用歸約推下界的想法是反證：假設 A 能做得比下界快，透過歸約讓 B 也變快，和 B 的下界矛盾
（資結書第~17.6.7~節用排序推出凸包的 $\Omega(n \lg n)$ 就是這樣）。 套到本題：假設 A 能在
$T(n) = o(n \lg n)$ 時間內解，題目給的事實只保證 B 能在
\[
O(n \log n + T(n)) = O(n \log n)
\]
時間內解。 這和 B 的下界 $\Omega(n \log n)$ 完全相容，推不出矛盾。 問題出在歸約本身就要
$O(n \log n)$：B 的困難可能全部花在歸約的那 $n \log n$ 上，跟 A 沒有關係。 要能推出 A 的下界，
歸約的額外時間必須是 $o(n \log n)$，例如 $O(n)$。

具體的反例：令 B 是比較模型下的排序，有 $\Omega(n \log n)$ 的下界（資結書定理~15.1）；令 A 是
「找 $n$ 個數的最小值」。 題目的事實成立：不管 A 用多少時間 $T(n)$，B 都可以直接用合併排序在
$O(n \log n) \subseteq O(n \log n + T(n))$ 時間內解。 B 也確實有 $\Omega(n \log n)$ 的下界。 但 A
從頭到尾掃一遍就好，$O(n)$，沒有 $\Omega(n \log n)$ 的下界。 所以一般情況下不能下這個結論。
\end{pastsoln}
